\section{Заключение}

В данной работе рассматривалось влияние LLM/FinBERT-generated features на прогнозирование финансового временного ряда. В отличие от постановки, где основной целью является поиск модели с максимальной точностью, работа была построена вокруг анализа признаков. Основной исследовательский вопрос состоял в том, дают ли признаки, извлечённые из финансовых текстов с помощью FinBERT, самостоятельный сигнал и добавочную предсказательную силу относительно стандартных рыночных признаков.

В теоретической части были рассмотрены три группы работ. Первая связана с прогнозированием финансовых временных рядов и ролью информации в динамике цен \cite{fama1970efficient}. Вторая показывает, что финансовые тексты, новости и социальные медиа могут быть источником признаков, отражающих ожидания и настроения участников рынка \cite{tetlock2007media,loughran2011liability,kearney2014textual,bollen2011twitter}. Третья связана с развитием языковых моделей: Transformer и BERT сделали возможным получение контекстных текстовых представлений, а FinBERT адаптирует этот подход к финансовому языку \cite{vaswani2017attention,devlin2019bert,araci2019finbert,yang2020finbert}. На основе этого был обоснован feature-centric подход: языковая модель используется не как самостоятельный прогнозист, а как генератор признаков для downstream-модели.

Эмпирическая часть была сформулирована как задача бинарной классификации направления цены закрытия акции AAPL на следующий торговый день:

\[
y_t = \mathbf{1}\{Close_{t+1} > Close_t\}.
\]

Для проверки влияния признаков были сформированы три feature groups:

\[
X_{price} = \{p_t\}, \qquad X_{text} = \{e_t\}, \qquad X_{combined} = \{[p_t;e_t]\},
\]

где $p_t \in \mathbb{R}^{11}$ --- рыночные признаки, $e_t \in \mathbb{R}^{768}$ --- FinBERT embedding, а $[p_t;e_t] \in \mathbb{R}^{779}$ --- объединённый вектор признаков. Такая структура позволила проверить не только абсолютное качество моделей, но и вклад каждой группы признаков.

Основной эксперимент RUN 03 был построен с контролем утечки данных: использовалось хронологическое разбиение train/validation/test без пересечения дат, а нормализация рыночных признаков обучалась только на train subset. Такой протокол важен для финансовых временных рядов, поскольку случайное перемешивание наблюдений может завышать качество модели и создавать иллюзию прогнозного сигнала.

Результаты показали, что модель \texttt{FinBERT-only} получила accuracy 0.550 против 0.532 у majority baseline. Разность составляет

\[
\Delta_{FinBERT-majority}^{Accuracy} = 0.550 - 0.532 = 0.018.
\]

Это указывает на наличие слабого самостоятельного текстового сигнала. Однако ROC-AUC для \texttt{FinBERT-only} равен 0.512, то есть остаётся близким к случайному уровню. Поэтому этот результат нельзя интерпретировать как доказательство устойчивой способности FinBERT-признаков прогнозировать направление цены.

Сравнение \texttt{prices-only} и \texttt{prices + FinBERT} показало положительную добавку от текстовых признаков относительно рыночной модели:

\[
\Delta_{text}^{Accuracy} = M_{Accuracy}(X_{combined}) - M_{Accuracy}(X_{price}) = 0.486 - 0.410 = 0.076.
\]

Аналогично для balanced accuracy:

\[
\Delta_{text}^{Balanced\ Accuracy} = 0.474 - 0.400 = 0.074.
\]

Следовательно, добавление FinBERT embeddings улучшает слабую модель на рыночных признаках. Однако этот эффект не является устойчивым в более строгом сравнении: объединённая модель не превосходит majority baseline по accuracy и уступает \texttt{FinBERT-only}. Поэтому комплементарность рыночных и текстовых признаков в текущем протоколе не подтверждается.

Дополнительная проверка с LSTM также не дала преимущества. Несмотря на то что LSTM теоретически может учитывать временную структуру признаков \cite{hochreiter1997lstm}, в данном эксперименте она не улучшила результат относительно более простой логистической регрессии. Это подтверждает, что ключевым ограничением является не только архитектура модели, но и информативность признакового пространства.

Главный вывод работы можно сформулировать следующим образом. В выбранном MVP-протоколе FinBERT-generated features дают слабый самостоятельный сигнал и улучшают результат относительно \texttt{prices-only}, но не демонстрируют устойчивой добавочной предсказательной силы относительно всех baseline-сравнений. Следовательно, корректно говорить не о неэффективности LLM-признаков в целом, а о том, что их влияние оказалось ограниченным в рамках выбранного актива, периода, горизонта прогноза и способа построения embedding-признаков.

Полученный результат согласуется с современной литературой по текстовому и мультимодальному финансовому прогнозированию: текстовые признаки могут быть полезны, но их эффект зависит от источника данных, временного выравнивания, таргета, горизонта прогноза и качества экспериментального протокола \cite{hiew2019bert,zou2022prebit,jiang2023finbert,albada2025bertweet,vuong2025llm}. Поэтому дальнейшая работа должна быть направлена не только на усложнение моделей, но и на более точную проверку признаков.

К возможным направлениям продолжения исследования относятся:

\begin{enumerate}
    \item расширение эксперимента на несколько активов и разные рыночные режимы;
    \item проверка других горизонтов прогноза, включая внутридневные и многодневные горизонты;
    \item сравнение FinBERT embeddings с sentiment-score, topic features и emotion features;
    \item использование rolling-window validation вместо одного фиксированного split;
    \item проверка fine-tuning языковой модели под конкретную финансовую задачу;
    \item анализ статистической устойчивости эффекта признаков на нескольких временных периодах.
\end{enumerate}

Таким образом, работа показывает, что LLM/FinBERT-generated features являются теоретически обоснованным источником признаков для финансового прогнозирования, но их полезность должна проверяться через строгие baseline-сравнения. В рамках проведённого эксперимента текстовые признаки не дали убедительного устойчивого преимущества, однако сама feature-centric методология позволяет аккуратно оценивать их вклад и может быть расширена в последующих исследованиях.
